% Fichier généré automatiquement par tools/generate_corrections.py.
% Source : CIN/CIN-03-Transmetteurs/26_RoueMotrice/26_RoueMotrice.tex
% Statut : partiel (4/5 question(s) renseignée(s), 0 reprise(s)).
\begin{corrigebox}[Corrigé partiel extrait de la banque]
\CorrigeQuestion{1}
\CorrigeACompleter{Aucun élément de réponse n'est présent dans la source d'origine pour cette question.}
\CorrigeQuestion{2}
\begin{marginfigure}
\includegraphics[width=.4\linewidth]{26_cor_01_BIS}
\end{marginfigure}
\CorrigeQuestion{3}
\begin{center}
\begin{tabular}{cccc}
\hline
Repère de  & Module  & Nombre & Diamètre primitif  \\
la roue & $m$ (mm) & de dents $Z$ & $D$ (mm) \\
\hline
27 & 1,5 &16 & 24\\ 
35 & 1,5 &84 & 126\\ 
5   &1,5 &14 & 21\\ 
11 & 1,5 & 56 & 84 \\ 
16 &  1,5&  75& 112,5\\ \hline
\end{tabular}
\end{center}
\CorrigeQuestion{4}
Voir figure précédente. Si le moteur tourne dans le sens positif, la roue tourne dans le sens négatif.
\CorrigeQuestion{5}
Le rapport de réduction de la transmission est le suivant : 
$k=\dfrac{Z_{27} Z_{5} Z_{11} }{Z_{35} Z_{11} Z_{16}} = \dfrac{16\cdot 14}{84\cdot 75} =0,0355 $.

La vitesse de rotation de la roue est donc de $\SI{53,33}{tr.min^{-1}}$ soit \SI{5,59}{rad.s^{-1}}. 
On en déduit la vitesse du véhicule : $5,59 \times 0,15 = \SI{0,84}{m.s^{-1}}\simeq \SI{3}{km.h^{-1}}$.
\end{corrigebox}
